\section{The adjugate scheme in characteristic two}
\label{sec:adjugate}

We now study the coefficient formula \eqref{eq:selector-coefficient} over
\[
 K=\overline{\mathbb F}_2.
\]
In this characteristic the permanent equals the determinant.  This lets
us associate a plane curve and a finite adjugate scheme to every matrix
\(X\).  Vanishing of the gradient will force this finite scheme to have
large length.

Fix an integer \(k\ge2\), pairwise distinct elements
\(\theta_1,\ldots,\theta_k\in K\), and
\(D=\operatorname{diag}(\theta_1,\ldots,\theta_k)\).  Set
\begin{equation}
 \label{eq:adjugate-parameters}
 q=\left\lfloor\frac{k-1}{3}\right\rfloor,\qquad
 r=k-2q,\qquad c=r-1,\qquad T=\binom{q+2}{2}.
\end{equation}
Thus
\begin{equation}
 \label{eq:adjugate-degrees}
 2q+c=k-1,\qquad 0\le q\le c.
\end{equation}
For a polynomial, brackets denote extraction of the indicated monomial
coefficient.  We use the homogeneous degree-\(r\) polynomial
\begin{equation}
 \label{eq:adjugate-coefficient-polynomial}
 G_D(X)=[s^qt^q]\per_k(X+sI_k+tD)
       =[s^qt^q]\det(X+sI_k+tD).
\end{equation}
Statements about its gradient below refer to the common zero locus of
all its \(k^2\) entrywise partial derivatives; no equation \(G_D=0\) is
included in that definition.

For each fixed \(X\in\Mat_k(K)\), put
\begin{equation}
 \label{eq:adjugate-curve}
 A=tD+uX,\qquad M=sI_k+A,\qquad
 F_X(s,t,u)=\det M,\qquad C_X=V(F_X)\subset\PP^2_K.
\end{equation}
Here \(F_X\) has degree \(k\), with coefficient \(1\) at \(s^k\).
We abbreviate \(C_X\) to \(C\) when \(X\) is fixed.

\begin{lemma}[Projection of the auxiliary curve]
 \label{lem:adjugate-projection}
For every \(X\in\Mat_k(K)\), the curve \(C_X\) is reduced and the map
\[
 \pi:C_X\longrightarrow\PP^1_K,\qquad [s:t:u]\longmapsto[t:u]
\]
is finite flat of degree \(k\) and is \'{e}tale over a nonempty open
subset of the base.  Every irreducible component of \(C_X\) meets the
line \(u=0\), and at each point of that intersection \(M\) has rank
\(k-1\).  Moreover, \(\pi^*\cO_{\PP^1}(1)=\cO_C(1)\), and
\begin{equation}
 \label{eq:adjugate-splitting}
 \pi_*\cO_C=\bigoplus_{j=0}^{k-1}\cO_{\PP^1}(-j),
 \qquad
 \mathcal E:=\pi_*\cO_C(k-1)
       =\bigoplus_{j=0}^{k-1}\cO_{\PP^1}(k-1-j).
\end{equation}
\end{lemma}

\begin{proof}
The restriction
\begin{equation}
 \label{eq:adjugate-fixed-fiber}
 F_X(s,t,0)=\prod_{i=1}^k(s+t\theta_i)
\end{equation}
is squarefree.  If a positive-degree irreducible homogeneous factor
\(H\) occurred at least twice in \(F_X\), its restriction to \(u=0\)
could not vanish identically: otherwise \(u\) would divide \(F_X\),
contrary to \eqref{eq:adjugate-fixed-fiber}.  That restriction is a
nonzero homogeneous polynomial of the same positive degree as \(H\).
Its square would divide \eqref{eq:adjugate-fixed-fiber}, a contradiction.
Thus \(C\) is reduced.  The restriction of each irreducible factor is
likewise a positive-degree binary form, so it has a projective zero
over \(K\).  Hence every component meets \(u=0\).

The points on that line are
\[
 P_i=[-\theta_i:1:0]\qquad(1\le i\le k).
\]
At \(P_i\), the matrix \(M\) is diagonal with exactly one zero entry.
It has rank \(k-1\), and its \(i\)-th diagonal adjugate entry is
\(\prod_{j\ne i}(\theta_j-\theta_i)\ne0\).

The only possible base point of the projection is \([1:0:0]\), which
does not belong to \(C\) because \(F_X(1,0,0)=1\).
On the base chart \(t\ne0\), use \(v=u/t\) and \(y=s/t\).  The inverse
image has coordinate algebra
\[
 K[v,y]\big/\det(yI_k+D+vX).
\]
Its defining polynomial is monic of degree \(k\) in \(y\); consequently
this algebra is free over \(K[v]\), with basis
\(1,y,\ldots,y^{k-1}\).  The chart \(u\ne0\) gives the same conclusion,
so \(\pi\) is finite flat of degree \(k\).
At \(v=0\) the monic polynomial has the distinct roots
\(-\theta_1,\ldots,-\theta_k\).
Its resultant with its \(y\)-derivative is nonzero there, and hence
nonzero on an open subset of the base.  Over that open subset the
finite flat map is \'{e}tale.  In particular its generic algebra is
a product of finite separable field extensions, also when \(k\) is even.

The sections \(t,u\) have no common zero on \(C\), so they identify
\(\pi^*\cO_{\PP^1}(1)\) with \(\cO_C(1)\).
The two monic-polynomial bases satisfy
\[
 (s/u)^j=(t/u)^j(s/t)^j\qquad(0\le j<k)
\]
on their overlap.  These are exactly the transition functions of
\(\cO_{\PP^1}(-j)\).  This proves the first splitting in
\eqref{eq:adjugate-splitting}; tensoring by \(\cO_{\PP^1}(k-1)\)
and using the projection formula proves the second.
\end{proof}

The fixed distinct-root fiber has established reducedness and a finite
flat, generically separable projection.  We next use the adjugate to turn a vanished
coefficient into a lower bound on the length of a finite scheme.

Each entry of \(\adj M\) is a homogeneous form of degree \(k-1\).
Define its ideal on \(C\) intrinsically by
\begin{equation}
 \label{eq:adjugate-ideal}
 J=\operatorname{im}\bigl(\cO_C(1-k)^{k^2}
                  \xrightarrow{\ \text{entries of }\adj M\ }\cO_C\bigr),
 \qquad Z=V_C(J).
\end{equation}
Thus a local trivialization of \(\cO_C(k-1)\) turns the adjugate
entries into generators of \(J\).  Changing the trivialization
multiplies all generators by a unit and leaves the ideal unchanged.
By Lemma~\ref{lem:adjugate-projection}, each component of \(C\)
contains a point at which an adjugate entry is nonzero.  The support
of \(\cO_C/J\) therefore contains no component of \(C\), so \(Z\)
is a finite scheme.  It may have nonreduced structure.

\begin{lemma}[Generation by adjugate entries]
 \label{lem:adjugate-generation}
Every homogeneous form of degree \(k-1\) that vanishes
scheme-theoretically on \(Z\) is a constant \(K\)-linear combination
of the entries of \(\adj M\).
\end{lemma}

\begin{proof}
All sheaves on the base in this proof are on \(\PP^1_K\).
The adjugate entries define a morphism
\[
 \psi:\Mat_k(\cO)\longrightarrow\mathcal E,\qquad
       Y\longmapsto\tr(Y\adj M).
\]
The trace pairs matrix entries in transposed order, so its image is
the \(\cO\)-span of all adjugate entries.  We first show that
\begin{equation}
 \label{eq:adjugate-image}
 \operatorname{im}\psi=\pi_*(J(k-1)).
\end{equation}
Work on either standard affine chart of the base and trivialize
\(\cO(1)\).  Write its coordinate ring as \(R\), the affine
\(s\)-coordinate as \(y\), and the corresponding matrix as \(A_0\).
The curve algebra is \(B=R[y]/(f)\), where
\(f=\det(yI_k+A_0)\).  Let \(Q\) be the \(R\)-span in \(B\) of the
entries of \(\adj(yI_k+A_0)\).
The adjugate identity gives, in \(B\),
\[
 y\adj(yI_k+A_0)=-A_0\adj(yI_k+A_0).
\]
Thus \(Q\) is stable under multiplication by \(y\).  Since \(B\) is
generated by \(y\) over \(R\), \(Q\) is already the \(B\)-ideal
generated by these entries.  Restoring the local trivialization
proves \eqref{eq:adjugate-image} on both charts.

Put \(\mathcal K=\ker\psi\).  To lift every global section of
\(\operatorname{im}\psi\) to a constant matrix, it is enough to prove
\(H^1(\mathcal K)=0\).  We obtain this vanishing by comparing
\(\mathcal K\) with the image of the commutator with \(A\).

The matrix \(A=tD+uX\) has entries in \(\cO(1)\), so the commutator
defines the morphism
\[
 \chi:\Mat_k(\cO(-1))\longrightarrow\Mat_k(\cO),\qquad
       Y\longmapsto AY-YA.
\]
The adjugate commutes with \(A\), and cyclicity of trace gives
\(\psi\chi=0\).  Put \(\mathcal I=\operatorname{im}\chi\); then
\(\mathcal I\subset\mathcal K\).

At the base point \([t:u]=[1:0]\), the commutator is \([D,-]\).
On the matrix unit \(E_{ij}\) it acts by multiplication by
\(\theta_i-\theta_j\).  These multipliers are nonzero for \(i\ne j\),
including in characteristic two.  Hence the commutator has rank
\(k^2-k\) at that point.
At the same point, the images under \(\psi\) of the diagonal matrix
units are the classes of
\[
 h_i(y)=\prod_{j\ne i}(y+\theta_j)
 \quad\text{in}\quad K[y]\big/\prod_j(y+\theta_j).
\]
Evaluation at the \(k\) distinct roots gives a diagonal matrix with
nonzero diagonal, so these classes are independent and \(\psi\)
has rank \(k\), the rank of \(\mathcal E\).  Its generic rank is
therefore \(k\), and its generic kernel has dimension \(k^2-k\).
The rank of \(\chi\) at \([1:0]\) shows that its generic rank is at
least \(k^2-k\), while \(\mathcal I\subset\mathcal K\) gives the
opposite inequality.  Thus \(\mathcal I\) and \(\mathcal K\) have
equal generic rank.
The quotient \(\mathcal K/\mathcal I\) is therefore a coherent
torsion sheaf on \(\PP^1\).

We now use two exact sequences to compute \(H^1(\mathcal K)\).
The first is
\[
 0\longrightarrow\ker\chi\longrightarrow
 \cO(-1)^{k^2}\longrightarrow\mathcal I\longrightarrow0.
\]
Since \(H^1(\PP^1,\cO(-1))=0\) and coherent \(H^2\) vanishes on
\(\PP^1\), its cohomology sequence gives \(H^1(\mathcal I)=0\).
These two facts follow, for example, from the usual computation
with the two standard affine charts of \(\PP^1\).
Next, a coherent torsion sheaf on \(\PP^1\) has finite support and
vanishing \(H^1\).  The sequence
\[
 0\longrightarrow\mathcal I\longrightarrow\mathcal K
 \longrightarrow\mathcal K/\mathcal I\longrightarrow0
\]
therefore gives \(H^1(\mathcal K)=0\).
Finally, the sequence
\[
 0\longrightarrow\mathcal K\longrightarrow\cO^{k^2}
 \longrightarrow\operatorname{im}\psi\longrightarrow0
\]
shows that constant matrices surject onto
\[
 H^0(\PP^1,\operatorname{im}\psi)=H^0(C,J(k-1)),
\]
where the equality uses \eqref{eq:adjugate-image} and finite
pushforward.

Let \(P\) be a degree-\(k-1\) homogeneous form vanishing on the
scheme \(Z\).  Its restriction to \(C\) is a section of \(J(k-1)\),
so the preceding surjectivity expresses it on \(C\) as a constant
linear combination of the adjugate entries.
Their difference is a degree-\(k-1\) form vanishing on \(C\).
Since the homogeneous ideal of this plane hypersurface is generated
by \(F_X\), of degree \(k\), the difference is zero.
Thus the equality holds as an equality of homogeneous forms.
\end{proof}

\begin{proposition}[Length forced by the gradient]
 \label{prop:adjugate-length}
With the parameters in \eqref{eq:adjugate-parameters}, let
\(X\in\Mat_k(K)\) satisfy \(\nabla G_D(X)=0\).
Then its adjugate zero scheme \(Z\) satisfies
\[
 \length Z\ge T=\binom{q+2}{2}.
\]
\end{proposition}

\begin{proof}
Homogeneity and entrywise differentiation give
\[
 G_D(X)=[s^qt^qu^r]F_X,\qquad
 \frac{\partial F_X}{\partial X_{ij}}=u(\adj M)_{ji}.
\]
Thus the gradient hypothesis is precisely the entrywise condition
\begin{equation}
 \label{eq:adjugate-coefficient-zero}
 [s^qt^qu^c]\adj M=0,\qquad c=r-1.
\end{equation}
This calculation uses no division by the degree \(r\).
By Lemma~\ref{lem:adjugate-generation}, the functional
\(\Lambda=[s^qt^qu^c]\) consequently vanishes on every degree-\(k-1\)
form that vanishes scheme-theoretically on \(Z\).

The next comparison is the line-bundle case of Ga{\l}{\k a}zka's
catalecticant bound under this annihilation condition.  It uses the
multiplication map \(P\mapsto(Q\mapsto\Lambda(PQ))\), where \(P\) and
\(Q\) have degrees \(q\) and \(k-1-q\), respectively
\cite[Theorem~4 and Propositions~6, 7, and~9 in arXiv:1605.08005v1]{Galazka2017}.
The complementary monomials below compute this map directly over \(K\).

Suppose that a degree-\(q\) form
\[
 P=\sum_{i+j+\ell=q}a_{ij\ell}s^it^ju^\ell
\]
vanishes on \(Z\).  For each triple in this sum, the monomial
\[
 s^{q-i}t^{q-j}u^{c-\ell}
\]
has nonnegative exponents because \(i,j,\ell\le q\le c\).
Its product with \(P\) still vanishes on \(Z\), and has degree
\(2q+c=k-1\).  Exactly one term contributes to its selected
coefficient, giving
\[
 0=\Lambda\bigl(s^{q-i}t^{q-j}u^{c-\ell}P\bigr)=a_{ij\ell}.
\]
Hence \(P=0\).  Restriction therefore induces an injection
\[
 H^0(\PP^2,\cO(q))\lhook\joinrel\longrightarrow
 H^0(Z,\cO_Z(q)).
\]
The space on the left has dimension \(\binom{q+2}{2}\).
On a finite scheme over \(K\), an invertible sheaf is free of rank
one on each Artin local component, so its space of global sections
has dimension equal to the scheme's length.  The space on the right
thus has dimension \(\length Z\), including all nonreduced
multiplicities.  The stated inequality follows.
\end{proof}
